% !TeX root = ../../Book2.tex
% 纲要标题：### 7.1 投射模
% 纲要位置：计划纲要.md，第 1274 行。

\section{投射模}
\label{b2:sec:0701}

若一个模有基，从它出发的同态便能沿满射提升：先为基向量的像选取原像，再线性延伸。没有基的模也可能具有这项性质。本节把它单独作为定义，并说明它恰好描述自由模的直和因子。除特别说明外，环均含幺元，模均为幺性左模。

% 纲要要点（供目录核对）：
%
% * 提升性质与自由模直和因子
% * 投射模和分裂正合列
% * 有限生成投射模的例子与幂等自同态的像；以有限对偶基就地证明 Hom 的张量实现
%

\subsection{提升性质与自由模直和因子}
\label{b2:subsec:0701:01}

\begin{definition}\label{b2:def:projective-module}
设 \(R\) 是环，\(P\) 是左 \(R\) 模。若对任意左 \(R\) 模之间的满同态 \(q:B\to C\) 和同态 \(f:P\to C\)，都存在同态 \(g:P\to B\) 满足 \(qg=f\)，就称 \(P\) 为\term[toushemo]{投射模}[projective module]。这里要求提升存在，不要求它唯一。
\end{definition}

这里要求下图中的虚箭头存在，并使两个从 \(P\) 到 \(C\) 的复合相等：
\[
\begin{tikzcd}[column sep=3em,row sep=2.5em]
& B \arrow[d,two heads,"q"] \\
P \arrow[r,"f"'] \arrow[ur,dashed,"g"] & C
\end{tikzcd}
\]

已知 \(f\) 的像在下面的 \(C\) 中；提升 \(g\) 则把这份映射沿满射抬到上面的 \(B\)，再由 \(q\) 映回时恢复 \(f\)。\cref{b2:prop:free-lifting}已经证明自由模具有这项性质。零模也投射。作为一个相反的例子，\(\mathbb Z/n\mathbb Z\) 在 \(n>1\) 时不是投射 \(\mathbb Z\) 模：若恒等映射能沿 \(\mathbb Z\to\mathbb Z/n\mathbb Z\) 提升，则 \(\bar1\) 的像应是被 \(n\) 零化的整数，只能为零，却又必须映回 \(\bar1\)。

\begin{theorem}\label{b2:thm:projective-characterizations}
设 \(R\) 是环，\(P\) 是左 \(R\) 模。下列条件等价：
\begin{equivalences}
\item\label{b2:item:projective-lifting} \(P\) 是投射模。
\item\label{b2:item:projective-split-epimorphism} 每个以 \(P\) 为终点的满同态都有模同态截面。
\item\label{b2:item:projective-summand} 存在左模 \(Q\)，使 \(P\oplus Q\) 是自由模。
\end{equivalences}
\end{theorem}
\begin{proof}
若 \(P\) 投射，对满同态 \(q:B\to P\) 提升 \(\operatorname{id}_P\)，便得到 \(qs=\operatorname{id}_P\)，即一个截面。

任意模都有自由满射 \(\pi:F\to P\)，例如取以 \(P\) 的底集合为基的自由模，把对应于 \(p\) 的基向量送到 \(p\)。若这个满射有截面，则\cref{b2:thm:split-exact-sequence}给出 \(F\cong\ker\pi\oplus P\)，所以第二项蕴含第三项。

最后设 \(F=P\oplus Q\) 自由，记包含与投影为 \(i:P\to F\)、\(p:F\to P\)，因而 \(pi=\operatorname{id}_P\)。给定定义中的 \(q\) 和 \(f\)，先用\cref{b2:prop:free-lifting}提升 \(fp:F\to C\)，得到 \(h:F\to B\) 且 \(qh=fp\)。令 \(g=hi\)，则 \(qg=fpi=f\)。
\end{proof}

定义要处理从 \(P\) 出发的所有提升问题，第二个条件把它缩成了以 \(P\) 为商模的分裂问题：每个短正合列 \(0\to A\to B\to P\to0\) 都应分裂。还可以只固定一个自由满射 \(\pi:F\to P\) 来检验；它有截面，就使 \(P\) 成为自由模的直和因子，因而投射。反过来，\(P\) 投射时，任何这样的自由满射都有截面。

\ProseParagraphBreak
这里“直和因子”包含嵌入与投影相容的意思。只有一个抽象单射 \(P\to F\)，并不足以证明 \(P\) 投射；\secref{b2:sec:0704}将给出嵌入自由模却不平坦、更不投射的理想。

\begin{proposition}\label{b2:prop:projective-sums-summands}
设 \(R\) 是环。左 \(R\) 投射模的任意直和是投射模，投射模的直和因子仍是投射模。
\end{proposition}
\begin{proof}
对 \(f:\bigoplus_iP_i\to C\) 和满射 \(q:B\to C\)，分别提升各限制 \(f_i:P_i\to C\)，得到 \(g_i:P_i\to B\)。\cref{b2:thm:module-sum-universal}将它们唯一合成 \(g:\bigoplus_iP_i\to B\)；\(qg=f\) 可在每个直和分量上验证。无限族的提升选择使用选择公理。

若 \(P\) 是投射模 \(T\) 的直和因子，用 \(i:P\to T\)、\(p:T\to P\) 且 \(pi=\operatorname{id}_P\) 表示它。先提升 \(fp\)，再与 \(i\) 复合，即得到 \(f\) 的提升。
\end{proof}

\subsection{投射模和分裂正合列}
\label{b2:subsec:0701:02}

投射性也可以用正合性表达。下面第一列的左端单射来自 \(A\to B\) 的单射，中间位置的正合性则说明：一个 \(P\to B\) 的同态落入 \(A\) 在 \(B\) 中的像，恰好等价于它复合到 \(C\) 为零。尚未保证的是每个 \(P\to C\) 都能提升到 \(B\)，这正是投射性要补足的末端满射。第二列将来源和目标的角色对换，下一节会用它把内射性解释为同态的延拓。

\begin{proposition}[Hom 的左正合性]\label{b2:prop:hom-left-exact}
设 \(R\) 是环，\(0\to A\xrightarrow{i}B\xrightarrow{q}C\to0\) 是左 \(R\) 模的短正合列，\(P,E\) 是左 \(R\) 模。下面两列交换群正合：
\[
0\longrightarrow\operatorname{Hom}_R(P,A)
\xrightarrow{i\circ-}\operatorname{Hom}_R(P,B)
\xrightarrow{q\circ-}\operatorname{Hom}_R(P,C),
\]
\[
0\longrightarrow\operatorname{Hom}_R(C,E)
\xrightarrow{-\circ q}\operatorname{Hom}_R(B,E)
\xrightarrow{-\circ i}\operatorname{Hom}_R(A,E).
\]
\end{proposition}
\begin{proof}
第一列中，\(i\) 单射保证 \(i\circ-\) 单射。若 \(qf=0\)，则 \(f(P)\subseteq\ker q=i(A)\)，由\cref{b2:prop:short-exact-kernel-cokernel}，\(f\) 唯一写成 \(ih\)，其中 \(h:P\to A\) 是同态。这给出中间位置的正合性。

第二列中，\(q\) 满射保证 \(-\circ q\) 单射。若 \(hi=0\)，则 \(h\) 在 \(i(A)\) 上为零，故公式 \(\bar h\bigl(q(b)\bigr)=h(b)\) 与代表元无关，并定义 \(\bar h:C\to E\)。它满足 \(h=\bar hq\)，且由 \(q\) 满射唯一确定。两列相邻映射的复合均为零，因此这些核与像的等式就是所需正合性。
\end{proof}

\begin{corollary}\label{b2:cor:projective-hom-exact}
设 \(R\) 是环，\(P\) 是左 \(R\) 模。\(P\) 投射，当且仅当函子 \(\operatorname{Hom}_R(P,-)\) 把每个左 \(R\) 模的短正合列变成交换群的短正合列；也等价于每个左 \(R\) 模短正合列 \(0\to A\to B\to P\to0\) 都分裂。
\end{corollary}
\begin{proof}
由\cref{b2:prop:hom-left-exact}，只需补上最后一个 Hom 映射的满射性，而它正是每个 \(P\to C\) 能沿 \(B\to C\) 提升的条件。任一满射与其核组成短正合列，所以这个条件覆盖定义中的所有满射。第二个等价性是\cref{b2:thm:projective-characterizations}与\cref{b2:thm:split-exact-sequence}的合用：满射有截面恰好使相应短正合列分裂。
\end{proof}

这个结论只要求右端的 \(P\) 投射。若 \(P\) 出现在左端，短正合列未必分裂。例如 \(0\to\mathbb Z\xrightarrow{\times n}\mathbb Z\to\mathbb Z/n\mathbb Z\to0\) 的左端自由，\(n>1\) 时仍不分裂。

\subsection{有限生成投射模的例子}
\label{b2:subsec:0701:03}

\begin{proposition}\label{b2:prop:finite-projective-summand}
设 \(R\) 是环，\(P\) 是左 \(R\) 模。\(P\) 有限生成且投射，当且仅当它是某个有限秩自由左 \(R\) 模 \(R^n\) 的直和因子；也等价于存在 \(n\ge0\) 和幂等自同态 \(E\in\operatorname{End}_R(R^n)\)，使 \(P\cong\operatorname{im}E\)。若 \(R\) 交换，后一条件可以用标准基写成幂等矩阵 \(E\in M_n(R)\)、\(E^2=E\) 的像。
\end{proposition}
\begin{proof}
若 \(P\) 有限生成，取满射 \(q:R^n\to P\)；若又投射，该满射有截面 \(s:P\to R^n\)，因而由\cref{b2:thm:split-exact-sequence}得到直和分解。反过来，有限自由模的直和因子由\cref{b2:prop:projective-sums-summands}知是投射模，而且是 \(R^n\) 在投影下的像，由标准基的有限个像生成。

任一这样的直和分解都有包含 \(s:P\to R^n\) 与投影 \(q:R^n\to P\)，满足 \(qs=\operatorname{id}_P\)。于是 \(E=sq\) 满足 \(E^2=s(qs)q=E\)，且 \(\operatorname{im}E=s(P)\cong P\)。反过来，若 \(E^2=E\)，每个 \(x\in R^n\) 都有分解
\[
x=E(x)+\bigl(x-E(x)\bigr),\qquad
E(x)\in\operatorname{im}E,\quad x-E(x)\in\ker E.
\]
若 \(y=E(x)\) 也在 \(\ker E\) 中，则 \(y=E^2(x)=E(y)=0\)，所以 \(R^n=\operatorname{im}E\oplus\ker E\)。因此 \(\operatorname{im}E\) 是有限自由模的直和因子。交换环上，标准基把自同态表示为通常作用于坐标列的矩阵，复合表示为矩阵乘法，故 \(E^2=E\) 正是幂等矩阵条件。
\end{proof}

设 \(e\in R\) 满足 \(e^2=e\)。作为左模，
\[
R=Re\oplus R(1-e).
\]
每个 \(r\) 都是 \(re+r(1-e)\)；若 \(x\) 同时属于两项，则右乘 \(e\) 分别得到 \(xe=x\) 与 \(xe=0\)，所以 \(x=0\)。因此 \(Re\) 是有限生成投射模。这一例子不要求 \(e\) 位于中心；右乘 \(e\) 才是这里的左模投影。

\ProseParagraphBreak
取 \(R=k\times k\)、\(e=(1,0)\)，则 \(Re=k\times0\) 非零，而 \((0,1)\) 零化的是整个模 \(Re\)，所以 \(\operatorname{Ann}_R(Re)\ne0\)。任意非零自由模却是忠实的：若 \(r\) 零化整个自由模，取其中一个基向量 \(e_i\)，则 \(re_i=0\) 的唯一坐标迫使 \(r=0\)。因此这个投射模不自由。这里必须比较整个模的零化子；仅有一个非零元素被非零标量零化，并不能排除自由性，正则自由模 \(R\) 中的 \(e\) 就被 \((0,1)\) 零化。

\ProseParagraphBreak
另一例是 \(R=M_n(k)\)、\(e=E_{11}\)：\(Re\) 由仅第一列可能非零的矩阵组成，可识别为通常的列向量模 \(k^n\)。当 \(n>1\) 时它也不自由，因为其 \(k\) 维数为 \(n\)，而非零自由 \(R\) 模若有有限基，其 \(k\) 维数为 \(n^2\) 的正整数倍。若有无限基，就可将它识别为 \(R^{(\Lambda)}\)，其中 \(\Lambda\) 无限。任意有限多个不同的标准基向量在 \(k\) 上也线性无关：在第 \(\lambda\) 个坐标读取一个线性关系，得到 \(a_\lambda I_n=0\)，故 \(a_\lambda=0\)。于是底层 \(k\) 空间的维数无限，与 \(\dim_k Re=n\) 矛盾。

\begin{proposition}[有限对偶基判据]\label{b2:prop:finite-dual-basis}
设 \(R\) 是环，\(P\) 是左 \(R\) 模。\(P\) 有限生成且投射，当且仅当存在有限个 \(p_1,\ldots,p_n\in P\) 和左 \(R\) 模同态 \(\varphi_i:P\to {}_RR\)，使对每个 \(p\in P\) 都有
\[
p=\sum_{i=1}^n\varphi_i(p)p_i.
\]
这样的元素与同态称为一组\term[youxian-duiouji]{有限对偶基}[finite dual basis]。它不要求 \((p_i)\) 线性无关，不要求 \(\varphi_i(p_j)=\delta_{ij}1_R\)，也不要求元素的系数表达唯一。
\end{proposition}
\begin{proof}
由\cref{b2:prop:finite-projective-summand}，有限生成投射模有分裂满射 \(q:R^n\to P\) 和截面 \(s:P\to R^n\)。令 \(p_i=q(e_i)\)，记第 \(i\) 个左坐标映射为 \(\pi_i:R^n\to R\)，并令 \(\varphi_i=\pi_i s\)。也就是说，先把 \(p\) 嵌回有限自由模，再读出第 \(i\) 个坐标。于是
\[
p=qs(p)=q\left(\sum_i\varphi_i(p)e_i\right)
=\sum_i\varphi_i(p)p_i.
\]

反过来，用 \(q(e_i)=p_i\) 定义 \(q:R^n\to P\)，用 \(s(p)=\sum_i\varphi_i(p)e_i\) 定义 \(s:P\to R^n\)。由于 \(\varphi_i(rp)=r\cdot\varphi_i(p)\)，\(s\) 确为左模同态。给定等式说明 \(qs=\operatorname{id}_P\)，所以 \(P\) 是有限自由模的直和因子。
\end{proof}

令 \(\psi_i:R\to P\)、\(r\mapsto rp_i\)，则重构公式也写成 \(\operatorname{id}_P=\sum_i\psi_i\varphi_i\)。每个复合都经过正则模 \(R\)，其像包含在循环子模 \(Rp_i\) 中；\(\varphi_i\) 和 \(\psi_i\) 都是左 \(R\) 线性的，因此这个解释在非交换环上也成立。有限对偶基是从有限自由模的直和分解继承的有限重构办法，允许 \((p_i)\) 有冗余。

\ProseParagraphBreak
沿用 \(R=k\times k\)、\(e=(1,0)\) 与 \(P=Re\) 的例子，取 \(p_1=e\)，并取包含映射 \(\varphi_1:P\hookrightarrow R\)。对 \(p=(a,0)\in P\)，有
\[
p=\varphi_1(p)e=(a,0)e,\qquad
\varphi_1(e)=e\ne1_R=(1,1).
\]
因此这一有限对偶基并不满足普通对偶基的 \(\varphi_i(p_j)=\delta_{ij}1_R\)。表示 \(p\) 的系数也不唯一：对每个 \(b\in k\)，\((a,b)e=(a,0)\)。有限对偶基给出一套指定的线性重构办法，却未必提供无关系的唯一坐标。

\ProseParagraphBreak
有限对偶基还把从 \(P\) 出发的 Hom 写成张量积。\cref{b2:prop:finite-free-hom-tensor}在有限自由情形用基与坐标泛函，把同态记录为有限多个基向量的像。对有限生成投射模，虽然 \((p_i)\) 不一定是基，有限对偶基仍使每个同态由有限多个值 \(f(p_i)\) 重构。正向对应仍把泛函 \(\lambda\) 与元素 \(m\) 组装为 \(p\mapsto\lambda(p)m\)；为了保留非交换环上的侧别，把额外的右作用一起写入陈述。

\begin{proposition}\label{b2:prop:finite-projective-hom-tensor}
设 \(R,S\) 为环，\({}_RP_S\) 为 \((R,S)\) 双模，且 \(P\) 作为左 \(R\) 模有限生成并投射。令 \(P^*=\operatorname{Hom}_R(P,{}_RR)\)，其 \((S,R)\) 双模作用为
\[
(s\lambda)(p)=\lambda(ps),\qquad
(\lambda r)(p)=\lambda(p)r.
\]
对每个左 \(R\) 模 \(M\)，令 \(F(M)=\operatorname{Hom}_R(P,M)\)，左 \(S\) 作用由 \((sf)(p)=f(ps)\) 给出；对左 \(R\) 模同态 \(a:M\to M'\)，令 \(F(a)(f)=af\)。则有关于 \(M\) 自然的左 \(S\) 模同构
\[
P^*\otimes_RM\longrightarrow F(M),\qquad
\lambda\otimes m\longmapsto\bigl(p\mapsto\lambda(p)m\bigr).
\]
\end{proposition}
\symidx{\(P^*=\operatorname{Hom}_R(P,R)\)：左模的右对偶模}
\begin{proof}
按式\eqref{b2:eq:hom-left-action}，\(P\) 的右 \(S\) 作用通过预复合，在 \(P^*\) 与 \(F(M)\) 上给出左 \(S\) 作用。\(P^*\) 的右 \(R\) 作用是在取值后右乘，也保留左 \(R\) 线性。两侧作用满足 \(\bigl((s\lambda)r\bigr)(p)=\lambda(ps)r=\bigl(s(\lambda r)\bigr)(p)\)，故 \(P^*\) 确为双模。由 \(\lambda(rp)=r\lambda(p)\)，\(p\mapsto\lambda(p)m\) 为左 \(R\) 线性。所写两变量映射双可加，且 \((\lambda r)(p)m=\lambda(p)(rm)\)，所以对 \(R\) 平衡；它也保留左 \(S\) 作用，因两种作用路径都在 \(p\) 处给出 \(\lambda(ps)m\)。故诱导所需同态。

由\cref{b2:prop:finite-dual-basis}，取有限个 \(p_i\in P,\lambda_i\in P^*\)，使 \(p=\sum_i\lambda_i(p)p_i\)。对任意同态 \(f:P\to M\)，左 \(R\) 线性给出 \(f(p)=\sum_i\lambda_i(p)f(p_i)\)。因此应把这些有限个值 \(f(p_i)\) 分别与重构泛函 \(\lambda_i\) 配对，反向定义 \(f\mapsto\sum_i\lambda_i\otimes f(p_i)\)。正向复合在 \(p\) 处取值为 \(\sum_i\lambda_i(p)f(p_i)=f(p)\)。另一复合将 \(\lambda\otimes m\) 送到
\[
\sum_i\lambda_i\otimes\lambda(p_i)m
=\left(\sum_i\lambda_i\cdot\lambda(p_i)\right)\otimes m
=\lambda\otimes m.
\]
最后一个等式来自对偶基公式施加 \(\lambda\)：\(\lambda(p)=\sum_i\lambda_i(p)\lambda(p_i)\)，其系数次序如式所写。故两映射互逆；正向映射左 \(S\) 线性，其逆也左 \(S\) 线性。对 \(M\) 上的同态，两种路径都将 \(\lambda\otimes m\) 送到 \(p\mapsto\lambda(p)a(m)\)，证明自然性。
\end{proof}

这个同构说明，有限生成投射模虽然未必有唯一的基坐标，有限对偶基仍足以把同态重构为有限张量和。在非交换环上，\(\lambda(p)\) 作用在 \(m\) 左侧，\(\lambda_i\cdot\lambda(p_i)\) 则使用 \(P^*\) 的右 \(R\) 作用；这些次序不能互换。“对偶基”这一名称并不允许忽略标量的作用方向。
